\begin{lemma}
For each fixed \(0<\eta<1/2\), put \(c=\frac12-\eta\).  Then
\[
c\,\noderef{RamseyScale}(k)
\]
tends to infinity with \(k\).  Moreover,
\[
\frac{\log\bigl(c\,\noderef{RamseyScale}(k)\bigr)}{\log k}\to 2,
\]
equivalently
\[
\log\bigl(c\,\noderef{RamseyScale}(k)\bigr)=2\log k+o(\log k).
\]
\end{lemma}
\begin{proof}
Let \(c=\frac12-\eta\).  Since \(0<\eta<1/2\), we have \(c>0\).
By \(\noderef{RamseyScale}\),
\[
c\,\noderef{RamseyScale}(k)
=c\,\frac{k^2}{\log k}.
\]
We first prove that this tends to infinity.  Let \(B\in\mathbb R\) be an
arbitrary target lower bound.  For all sufficiently large \(k\) we have
\(k>1\), hence \(\log k>0\), and also the elementary estimate
\(\log k\le k\).  If \(B\le0\), then
\(c k^2/\log k>0\ge B\) for all such \(k\).  If \(B>0\), choose \(K\) so large
that \(K>1\) and \(cK\ge B\).  Then for every \(k\ge K\),
\[
c\,\frac{k^2}{\log k}\ge c\,\frac{k^2}{k}=ck\ge cK\ge B.
\]
Thus for every real \(B\), \(c\,\noderef{RamseyScale}(k)\ge B\) eventually,
which is exactly
\[
c\,\noderef{RamseyScale}(k)\to\infty.
\]

It remains to compute the logarithmic ratio.  The previous positivity also
shows that for all sufficiently large \(k\), the logarithms below are taken of
positive quantities.  Using \(c>0\), \(\log k>0\), and
\(\noderef{RamseyScale}(k)=k^2/\log k\), we have
\[
\begin{aligned}
\log\bigl(c\,\noderef{RamseyScale}(k)\bigr)
&=\log c+\log\left(\frac{k^2}{\log k}\right)\\
&=\log c+\log(k^2)-\log(\log k)\\
&=\log c+2\log k-\log\log k.
\end{aligned}
\]
Dividing by \(\log k\), which is positive eventually, gives
\[
\frac{\log\bigl(c\,\noderef{RamseyScale}(k)\bigr)}{\log k}
=2+\frac{\log c}{\log k}
 -\frac{\log\log k}{\log k}.
\]
The middle term tends to \(0\), because \(\log c\) is a fixed real number and
\(\log k\to\infty\).

We also have
\[
\frac{\log\log k}{\log k}\to0.
\]
Indeed, set \(t=\log k\).  Then \(t\to\infty\).  For \(t\ge1\), the elementary
calculus estimate \(\log t\le\sqrt t\) holds: the function
\(\sqrt t-\log t\) decreases on \(1\le t\le4\), remains positive at \(t=4\),
and then increases for \(t\ge4\).  Therefore, once \(t=\log k\ge1\),
\[
0\le\frac{\log\log k}{\log k}
=\frac{\log t}{t}
\le \frac{\sqrt t}{t}
=\frac1{\sqrt t},
\]
and the final expression tends to \(0\) because \(t\to\infty\).  Hence both
error terms in the displayed ratio tend to \(0\), so
\[
\frac{\log\bigl(c\,\noderef{RamseyScale}(k)\bigr)}{\log k}\to2.
\]
\end{proof}
